Fix an accepted finite fixing trace \(\tau\), fix a domain label \(s\) occurring in the trace, and fix \(\mathbf M\in\mathcal C(\mathcal P)\). By \(\mathrm{def:compatibility\text{-}fiber}\), the family \(\mathbf M\) satisfies \(\mathrm{ass:recursive\text{-}compatibility}\) and \(\mathrm{ass:district\text{-}factorization}\). Acyclicity and the finite response carrier of \(\mathrm{def:district\text{-}response\text{-}profile}\) imply that recursive substitution assigns to every event \(A\) occurring in the trace a total deterministic indicator
\[
 I_A:\prod_{D\in\mathcal D(G)}\Omega_D\longrightarrow\{0,1\}.
\tag{1}
\]
The vertex set, its district partition, every declared state space, and every response-profile set are finite; hence the product and every sum below are finite. Independence of the district response blocks, which is \(\mathrm{ass:district\text{-}factorization}\), gives
\[
 P_{M^s}(A)
 =\sum_{\omega\in\prod_{D\in\mathcal D(G)}\Omega_D}
 I_A(\omega)\prod_{D\in\mathcal D(G)}q_D^s(\omega_D).
\tag{2}
\]
Thus each unconditional response-profile event probability is a polynomial with integer coefficients in the finitely many coordinates \(q_D^s(\omega_D)\). The same argument applies separately to every source-labelled probability leaf. Algebraic combinations of leaves are therefore polynomials in the union of their coordinate families; this observation does not posit any joint probability law across distinct domain labels.

We first justify every strict-positivity certificate that can guard a divisor in \(\tau\). Let \(x\) collect the finitely many response-law coordinates named by such a certificate. The typing rules of \(\mathrm{def:finite\text{-}extraction\text{-}calculus}\) require its declared compatibility fiber to be represented by checked constraints
\[
 h_i(x)=0\quad(1\le i\le r),
 \qquad
 g_j(x)\ge 0\quad(1\le j\le t),
\tag{3}
\]
and require the response coordinates of every model in that fiber to satisfy these constraints. For a rational-constant certificate, exact normalization verifies a rational lower bound \(\varepsilon>0\) for the certified polynomial \(p\). For a linear face certificate, the checker verifies rational multipliers and the polynomial identity
\[
 p(x)-\varepsilon
 =\sum_{i=1}^{r}\alpha_i h_i(x)
  +\sum_{j=1}^{t}\lambda_j g_j(x),
 \qquad \varepsilon>0,\quad \lambda_j\ge 0.
\tag{4}
\]
For a preordering certificate, it verifies
\[
 p(x)-\varepsilon
 =\sum_{i=1}^{r}\tau_i(x)h_i(x)
  +\sum_{J\subseteq\{1,\ldots,t\}}
    \sigma_J(x)\prod_{j\in J}g_j(x),
 \qquad \varepsilon>0,
\tag{5}
\]
where every displayed \(\sigma_J\) is exactly checked to be a finite sum of squares of rational polynomials. Substitution of (3) into (4) gives \(p(x)\ge\varepsilon>0\). In (5), the equality terms vanish, each factor \(g_j(x)\) is nonnegative, every finite product of those factors is nonnegative, and every \(\sigma_J(x)\) is nonnegative; hence again \(p(x)\ge\varepsilon>0\). The rational-constant case has the same conclusion by its exact comparison. These cases exhaust the positivity constructors in \(\mathrm{def:finite\text{-}extraction\text{-}calculus}\). Consequently every context probability used as a divisor by the accepted trace is strictly positive at \(\mathbf M\).

We next verify the probability constructors. The false and true events have indicators zero and one. For a finite pairwise-disjoint family \((A_j:j\in J)\), pointwise equality of indicators gives
\[
 I_{\bigsqcup_{j\in J}A_j}=\sum_{j\in J}I_{A_j},
\]
and substitution in (2) yields
\[
 P_{M^s}\!\left(\bigsqcup_{j\in J}A_j\right)
 =\sum_{j\in J}P_{M^s}(A_j).
\tag{6}
\]
Marginalization is (6) applied to the finite partition indexed by the marginalized variable's declared state space. Pointwise \(I_{A^{\mathsf c}}=1-I_A\), so normalization of (2) gives
\[
 P_{M^s}(A^{\mathsf c})=1-P_{M^s}(A).
\tag{7}
\]
An algebraic product constructor preserves equality of already interpreted scalar expressions. If a product is annotated as the probability of a district-split event, its annotation is valid because \(\mathrm{ass:district\text{-}factorization}\) makes the indicated district response blocks independent. A type-correct substitution is a homomorphism for rational constants, finite sums, differences, and products, and therefore preserves every checked polynomial identity.

For a conditional constructor with numerator event \(A\cap H\) and context event \(H\), the certificate argument gives
\[
 P_{M^s}(H)>0.
\tag{8}
\]
The quotient is therefore defined, and the definition of conditional probability gives
\[
 P_{M^s}(A\mid H)
 =\frac{P_{M^s}(A\cap H)}{P_{M^s}(H)}.
\tag{9}
\]
Equations (2), (6), (7), and (9), by structural induction over the finite trace syntax in \(\mathrm{def:finite\text{-}extraction\text{-}calculus}\), identify the semantic value of every subexpression of \(\tau\).

Let \(L_\tau\) and \(R_\tau\) be the resulting response-profile expressions at the input and output of the trace, and let \(H_1,\ldots,H_J\) be its distinct divisor expressions. Clearing denominators uses a multiplier
\[
 \Delta_\tau(x)=\prod_{j=1}^{J}H_j(x)^{k_j},
 \qquad k_j\in\mathbb N_{>0},
\tag{10}
\]
with the empty product interpreted as one. By (8),
\[
 \Delta_\tau(x(\mathbf M))>0.
\tag{11}
\]
Acceptance of \(\tau\), by the exact-checker clause of \(\mathrm{def:finite\text{-}extraction\text{-}calculus}\), means that rational normalization and coefficient comparison verify the formal polynomial identity
\[
 \Delta_\tau L_\tau-\Delta_\tau R_\tau=0.
\tag{12}
\]
All coordinate sets and sums are finite, so substitution of the response-law coordinates \(x(\mathbf M)\) into (12) is valid and yields
\[
 \Delta_\tau(x(\mathbf M))L_\tau(x(\mathbf M))
 =\Delta_\tau(x(\mathbf M))R_\tau(x(\mathbf M)).
\]
Division by the strictly positive quantity in (11) proves
\[
 L_\tau(x(\mathbf M))=R_\tau(x(\mathbf M)).
\]
Since \(\mathbf M\in\mathcal C(\mathcal P)\) was arbitrary, every accepted finite fixing trace preserves its represented event probability throughout the compatibility fiber, as claimed.